\chapter{Particiones}

    \section{INTRODUCCIÓN}

    Hasta ahora este libro ha tratado principalmente de la \textit{teoría multiplicativa de
    números}, un estudio de las funciones aritméticas relacionadas con la descomposición de
    enteros en factores primos. Ahora se vuelve hacia otra rama de la Teoría de números llamada
    \textit{teoría aditiva de números}. Un problema básico en este contexto consiste en
    expresar un entero positivo dado $n$ como suma de los enteros de un cierto conjunto $A$,
    dado de antemano,
    \begin{equation*}
        A=\{a_1,a_2,\dots\},
    \end{equation*}
    en donde los elementos $a_i$ son números especiales tales como primos, cuadrados, cubos,
    números triangulares, etc. Cada representación de $n$ como suma de elementos de $A$ se
    llama \textit{partición} de $n$ y nos interesa la función $A(n)$ que cuenta el número de
    particiones de $n$ en sumandos de $A$. Ilustraremos el problema con algunos ejemplos
    famosos.

    \textbf{Conjetura de Goldbach.} \textit{Todo número par $n>4$ es suma de dos primos
    impares.}

    En este ejemplo $A(n)$ es el número de soluciones de la ecuación
    \begin{equation}\label{eq:14-1}
        n=p_1+p_2,
    \end{equation}
    en donde los $p_i$ son primos impares. La afirmación de Goldbach es que $A(n)\geq1$ para
    todo par $n>4$. Esta conjetura data de 1742 y, en esta fecha, quedó sin decidir. En 1937 el
    matemático ruso Vinogradov probó que cada entero par suficientemente grande es la suma de
    tres primos impares. En 1966 el matemático chino Chen Jing-run demostró que cada número par
    suficientemente grande es la suma de un primo más un número que no posee más de dos
    divisores primos. (Ver \cite{chen1966}.)

    \textbf{Representación por cuadrados.} \textit{Para un entero dado $k\geq2$ consideramos la
    función de partición $r_k(n)$ que cuenta el número de soluciones de la ecuación}
    \begin{equation}\label{eq:14-2}
        n=x_1^2+\cdots+x_k^2,
    \end{equation}
    \textit{en donde los $x_i$ pueden ser positivos, negativos o nulos, y el orden de los
    sumandos se toma de acuerdo con el orden creciente de su valor.}

    Para $k=2,4,6$ ó $8$, Jacobi \cite{jacobi1829} expresó $r_k(n)$ en términos de las funciones
    divisor. Por ejemplo, demostró que
    \begin{equation*}
        r_2(n)=4\{d_1(n)-d_3(n)\},
    \end{equation*}
    en donde $d_1(n)$ y $d_3(n)$ son el número de divisores de $n$ congruentes con 1 y 3 mód 4,
    respectivamente. Entonces, $r_2(5)=8$ puesto que ambos divisores, 1 y 5, son congruentes
    con 1 mód 4. De hecho existen cuatro representaciones dadas por
    \begin{equation*}
        5=2^2+1^2=(-2)^2+1^2=(-2)^2+(-1)^2=2^2+(-1)^2,
    \end{equation*}
    y otras cuatro con el orden de los sumandos invertido.

    Para $k=4$ Jacobi probó que
    \begin{align*}
        r_4(n)&=8\sum_{d\mid n}d=8\sigma(n)&&\text{si }n\text{ es impar,}\\
        &=24\sum_{\substack{d\mid n\\d\text{ impar}}}d&&\text{si }n\text{ es par.}
    \end{align*}
    Las fórmulas para $r_6(n)$ y $r_8(n)$ son algo más complicadas pero del mismo tipo general.
    (Ver \cite{dickson1930}.)

    Se han obtenido hasta el momento fórmulas exactas para $k=3$, $5$, ó $7$; contienen la
    extensión de Jacobi del símbolo de Legendre para restos cuadráticos. Por ejemplo, si $n$ es
    impar sabemos que
    \begin{align*}
        r_3(n)&=24\sum_{m\leq n/4}(m\,|\,n)&&\text{si }n\equiv1\pmod{4}\\
        &=8\sum_{m\leq n/2}(m\,|\,n)&&\text{si }n\equiv3\pmod{4},
    \end{align*}
    en donde ahora los números $x_1$, $x_2$, $x_3$ de \eqref{eq:14-2} se han tomado primos entre
    sí.

    Para valores grandes de $k$, el análisis de $r_k(n)$ es considerablemente más complicado.
    Existe mucha literatura sobre este tema con contribuciones de Mordell, Hardy, Littlewood,
    Ramanujan, y muchos otros. Para $k\geq5$ sabemos que $r_k(n)$ se puede expresar mediante
    una fórmula asintótica de la forma
    \begin{equation}\label{eq:14-3}
        r_k(n)=\rho_k(n)+R_k(n),
    \end{equation}
    en donde $\rho_k(n)$ es el término principal, dado por la serie infinita
    \begin{equation*}
        \rho_k(n)=\frac{\pi^{k/2}n^{k/2-1}}{\Gamma\left(\dfrac{k}{2}\right)}\sum_{q=1}^{\infty}\sum_{\substack{h=1\\(h,q)=1}}^{q}\left(\frac{G(h;q)}{q}\right)^ke^{-2\pi inh/q},
    \end{equation*}
    y $R_k(n)$ es un resto de orden menor. La serie de $\rho_k(n)$ se llama \textit{serie
    singular} y los números $G(h;q)$ son las sumas cuadráticas de Gauss,
    \begin{equation*}
        G(h;q)=\sum_{r=1}^{q}e^{2\pi ihr^2/q}.
    \end{equation*}
    En 1917, Mordell observó que $r_k(n)$ es el coeficiente de $x^n$ en el desarrollo en serie
    de la $k$-ésima potencia de la serie
    \begin{equation*}
        \vartheta=1+2\sum_{n=1}^{\infty}x^{n^2}.
    \end{equation*}
    La función $\vartheta$ se halla relacionada con las funciones modulares elípticas que juegan
    un papel realmente importante en la obtención de \eqref{eq:14-3}.

    \textbf{Problema de Waring.} \textit{Determinar si para un entero positivo dado $k$, existe
    un entero $s$ (que dependa sólo de $k$) tal que la ecuación}
    \begin{equation}\label{eq:14-4}
        n=x_1^k+x_2^k+\cdots+x_s^k
    \end{equation}
    \textit{tenga soluciones para cada $n\geq1$.}

    El problema se conoce con el nombre del matemático inglés E. Waring que estableció en 1770
    (sin demostración y con limitada evidencia numérica) que cada $n$ es la suma de 4
    cuadrados, de 9 cubos, de 19 cuartas potencias, etc. En este ejemplo la función de partición
    $A(n)$ es el número de soluciones de \eqref{eq:14-4}, y el problema consiste en decidir si
    existe un número $s$ tal que $A(n)\geq1$ para todo $n$.

    Si $s$ existe para un $k$ dado entonces existe un mínimo valor de $s$ y éste se designa
    $g(k)$. Lagrange probó la existencia de $g(2)$ en 1770 y, en los 139 años siguientes, se
    probó la existencia de $g(k)$ para $k=3,4,5,6,7,8$ y $10$. En 1909, Hilbert demostró la
    existencia de $g(k)$ para todo $k$ mediante un argumento inductivo pero no determinó su
    valor numérico para cualquier $k$. Actualmente se conoce el valor exacto de $g(k)$ para
    cada $k$ excepto para $k=4$. Hardy y Littlewood hallaron una fórmula asintótica para el
    número de soluciones de \eqref{eq:14-4} en términos de una serie singular análoga a la de
    \eqref{eq:14-3}. Para una reseña histórica del problema de Waring, ver W. J. Ellison
    \cite{ellison1971}.

    \textbf{Particiones no restrictivas.} Uno de los problemas fundamentales en la teoría
    aditiva de números es el de las particiones no restrictivas. El conjunto de los sumandos
    consta de todos los enteros positivos, y la función de partición que se estudia es el número
    de formas en que podemos descomponer $n$ como suma de enteros positivos $\leq n$, esto es,
    el número de soluciones de
    \begin{equation}\label{eq:14-5}
        n=a_{i_1}+a_{i_2}+\cdots
    \end{equation}
    El número de sumandos no se halla sometido a ninguna restricción, se permite repetir
    sumandos, y el orden de los sumandos no se toma en cuenta. La función de partición
    correspondiente se designa $p(n)$ y se llama \textit{función de partición no restrictiva},
    o simplemente \textit{función de partición}. Los sumandos se llaman \textit{partes}. Por
    ejemplo, existen exactamente cinco particiones de 4, dadas por
    \begin{equation*}
        4=4=3+1=2+2=2+1+1=1+1+1+1,
    \end{equation*}
    luego $p(4)=5$.

